\chapter{绪\quad 论}

\section{研究背景与意义}

路网是城市交通系统的核心基础设施，承载着居民日常出行与物流运输的基本需求。随着城市化进程的持续推进，全球已有超过 68\% 的人口居住在城市地区\cite{UNDESA2018}，交通拥堵、安全和效率等问题日益严峻，智能交通系统（Intelligent Transportation Systems, ITS）的发展对提升交通运行效率具有重要意义。在交通预测领域，基于图神经网络的方法已成为主流技术路线——Li 等\cite{li2018dcrnn} 提出的扩散卷积循环神经网络（DCRNN）和 Yu 等\cite{yu2018stgcn} 提出的时空图卷积网络（STGCN）分别引领了图神经网络在交通流预测中的应用浪潮。Jin 等\cite{jin2024stgnn} 对这一领域进行了系统综述，梳理了近年来时空图神经网络在城市计算各应用场景中的发展脉络。路网表示学习（Road Network Representation Learning）作为智能交通的关键技术基础，旨在将路网中的拓扑结构和语义属性编码为低维稠密向量，以支撑路径规划\cite{wu2020hrnr}、行程时间估计\cite{chen2021toast}、平均速度推断\cite{wu2020hrnr}、轨迹相似性检索等核心下游任务。

真实路网区别于图结构的核心特征在于其显著的多尺度层次组织。Barth\'{e}lemy\cite{Barthelemy2011} 的研究表明，道路系统具有尺度不变的树状分支结构——少数动脉级主干道与大量细粒度地方支路构成典型的指数分支模式。这种层次性渗透在路网的多个维度中：在功能维度上，从高速公路、干道到居民区支路，路段按承载功能形成明确的层级体系；在结构维度上，地理相邻路段自然聚合为结构区域（如立交桥群、交叉口群落），进而组成功能分区（如商业区、居民区）；在拓扑维度上，Porta 等人\cite{PortaCrucitti2006a} 和 Crucitti 等人\cite{Crucitti2006} 证实路网的中心性度量呈幂律分布，少数关键节点对全网连通性具有主导作用。层次结构建模对路网分析任务的价值已得到实证验证：Wu 等人提出的层次路网表示（HRNR）\cite{wu2020hrnr} 通过构建功能区、结构区和路段三层神经架构，在行程时间估计和速度推断任务上取得了显著提升。此后，Toast\cite{chen2021toast}、START\cite{chen2022start}、JCLRNT\cite{mao2022jclrnt}、SARN\cite{li2023sarn} 等方法通过多样化的预训练目标和编码器设计进一步扩展了路网建模范式。

然而，上述方法均在欧氏空间中运行，对路网层次关系的建模依赖于赋值矩阵或注意力机制，而非在几何层面上利用层次结构。欧氏空间的体积随半径呈多项式增长（$V \propto r^d$），与树状层次结构中节点数量随深度指数增长的模式存在根本性不匹配，在低维空间中嵌入层次数据必然导致严重失真\cite{sarkar2011low,nickel2017poincare}。双曲几何以其恒定负曲率和指数体积增长特性，为层次结构数据的嵌入提供了天然的几何先验\cite{krioukov2010hyperbolic,peng2022hyperbolic}。Nickel 和 Kiela\cite{nickel2017poincare} 展示了庞加莱嵌入能够以显著低失真将 WordNet 名词层次嵌入仅 5 维的双曲空间；Sarkar\cite{sarkar2011low} 从理论上证明仅用 2 维双曲空间即可近乎无损地嵌入任意树结构。近年来，视觉-语言领域的 HyCoCLIP\cite{pal2025hycoclip} 进一步通过蕴含锥\cite{ganea2018entailment}机制在双曲空间中对图像"整体-局部"的层次关系施加了显式几何约束。路网中"功能区 $\supseteq$ 局部区域 $\supseteq$ 路段"的包含关系与此高度类似，这一结构类比为将双曲几何工具引入路网层次建模提供了坚实的语义基础。然而，目前尚无工作系统性地利用双曲空间的几何特性来建模路网的层次结构，如何将双曲几何的理论优势转化为路网表示学习的实际性能提升，仍是亟待解决的开放问题。

\section{研究问题与挑战}

欧氏空间对层次结构表示的存在局限。欧氏空间是当前深度学习中最常用的嵌入空间，但其在表示层次结构数据时存在本质局限。$d$ 维欧氏空间中，半径为 $r$ 的球内体积与 $r^d$ 成正比——这是多项式级的增长。然而，树状层次结构中节点数量随层级深度呈指数增长：在一棵分支因子为 $b$ 的平衡树中，第 $k$ 层有 $b^k$ 个节点。这种根本性的不匹配意味着，若将层次结构嵌入欧氏空间，要么需要极高维度，要么产生严重的失真\cite{sarkar2011low,nickel2017poincare}。Nickel 和 Kiela\cite{nickel2017poincare} 的研究表明，在低维欧氏空间中嵌入含 $n$ 个节点的树需要 $O(\log n)$ 维度，且存在无法消除的失真。

对路网而言，这一局限意味着：现有欧氏方法在建模层次关系时，缺乏几何层面的支持，只能通过软赋值矩阵或注意力机制间接表达，嵌入质量受限于空间的表达能力。

现有路网层次建模方法对层次关系的表达方式存在根本性的差距。在 HRNR 等方法中，层次关系通过可学习的赋值矩阵以软性方式表达——路段嵌入与其所属结构区域的嵌入之间，仅通过矩阵乘法和注意力机制关联，而非利用嵌入空间的几何结构显式约束"包含"关系。标准 GNN 的消息传递在平坦图上进行，远距离层次信息难以有效传播至叶节点层级，且存在过平滑（Over-smoothing）和过压缩（Over-squashing）问题\cite{topping2022oversquashing}。

这一差距引出了一个核心的方法论问题：能否找到一种嵌入空间，其几何特性天然适配路网的层次结构，使得层次关系的编码从依赖网络参数的隐式学习，转变为由空间几何支撑的显式约束。

跨城市场景下层次结构的泛化能力也是一个新的挑战。不同城市的路网呈现出迥异的形态特征：从具有高度不规则树状结构的中世纪有机布局（如威尼斯，$\delta=3.0$），到现代规划的网格状区域（如雅典，$\delta=32.0$）。现有路网表示方法通常在单一城市数据上训练，跨城市泛化能力有限，难以将在一个城市中学到的层次结构知识迁移到形态差异显著的新城市中。如果层次关系的编码依赖于特定城市的拓扑细节而非普遍的几何规律，那么跨城市泛化就会面临根本性困难。

\section{研究内容与贡献}


本文的核心思路是：以双曲几何为统一工具，探索路网层次结构建模。双曲空间的负曲率导致其体积随半径呈指数增长，与树状层次结构中节点数量随深度指数增长的规律天然契合。蕴含锥（Entailment Cones）\cite{ganea2018entailment}等双曲几何工具进一步提供了在嵌入空间中建模偏序关系的手段。

围绕这一思路，本文按照从隐式到显式的递进逻辑，开展两项互补的研究工作。

第一部分聚焦于路段级别的特征学习与分类任务，提出双曲路网图神经网络（Hyperbolic Road Graph Neural Network, HRGNN）。HRGNN 利用庞加莱球模型中的双曲几何，借助双曲空间的度量性质（双曲距离随层次深度增大而增大）隐式捕获路网的层次信息，同时通过双通道注意力机制平衡局部连通性与全局层次依赖的建模。这一工作验证了双曲空间对路网层次结构建模的基本有效性，同时也揭示了仅依赖空间度量性质进行隐式建模的局限——层次包含关系无法在嵌入空间中得到显式约束。

第二部分在 HRGNN 的基础上进一步推进，提出 HHRoad——一种在 Lorentz 双曲模型中运行的层次路网表示学习框架。HHRoad 使用组合蕴含学习的思想，通过蕴含锥机制对路网三层层次结构施加显式的几何约束，将层次包含关系从软赋值矩阵的隐式表达升级为双曲空间中蕴含锥不等式的硬几何约束，并结合层次对比学习捕获多尺度语义相似性。

HRGNN 回答了"双曲空间是否有助于路网层次建模"的基本问题，验证了双曲几何的有效性；HHRoad 则回答了"如何在双曲空间中显式约束路网层次关系"的深入问题，将层次建模从借助空间特性推进到设计几何约束。

本文的主要贡献可以总结为以下4个方面：

（1）\textbf{提出了首个专为交通路网设计的双曲图神经网络（HRGNN）}。通过将庞加莱球模型中的双曲运算引入路网表示学习，HRGNN 在保留层次结构几何先验的同时，实现了对路网多尺度特征的有效提取。双通道双曲注意力机制新颖地将局部邻域聚合与重要性引导 DFS 路径采样相结合，在双曲空间中同时建模路网的局部连通性和层次依赖关系。

（2）\textbf{提出了基于蕴含锥的组合蕴含学习范式用于路网}。将蕴含锥机制应用于路网三层层次结构，首次在路网表示学习中实现了对层次包含关系的显式几何约束，并设计了双曲对比学习框架，利用 Lorentz 距离捕获路段多尺度的语义相似性。

（3）\textbf{系统论证了路网层次结构与双曲空间的适配性}。通过覆盖 29 个国家 50 个城市的大规模 $\delta$-双曲性分析，从实证角度验证了路网拓扑的树状特性，并从理论角度论证了双曲空间作为路网层次建模几何框架的优势，提出了隐式捕获与显式约束两种范式的统一分析框架。

（4）\textbf{开展了系统性的实验评估}。在城市内分类、跨城市迁移学习和多城市多任务基准（VecCity）上对两项工作进行了全面评估，并通过消融实验、可视化分析和理论分析验证了各组件的贡献。

\section{论文组织结构}

本文共六章，组织结构如下：

第一章（本章）为绪论，介绍研究背景、动机、问题陈述、技术路线和主要贡献。

第二章为相关工作综述，系统梳理路网表示学习、图神经网络、双曲表示学习以及对比学习与蕴含学习的相关研究。

第三章为路网层次结构特性与双曲空间建模理论分析，对路网层次结构的多尺度特性进行系统分析，论证双曲空间与路网层次结构的适配性，并提出隐式捕获与显式约束两种建模范式，为后续两章方法提供统一的理论基础。

第四章为基于双通道注意力的双曲图神经网络（HRGNN），详细介绍 HRGNN 的模型设计与实验评估，展示隐式层次建模范式的有效性及其局限。

第五章为基于蕴含锥的双曲层次路网表示学习（HHRoad），详细介绍 HHRoad 的模型架构、蕴含学习机制、对比学习框架与实验结果分析，展示显式层次建模范式的优势。

第六章为总结与展望，综合总结前面的工作内容，总结理论启示，分析研究局限，并提出未来研究方向。